\documentclass[12px,a4paper]{article}
\usepackage[utf8]{inputenc}
\usepackage[spanish]{babel}
\usepackage{graphicx}
\usepackage{listings}
\usepackage{xcolor}
\usepackage{hyperref}
\usepackage{geometry}
\geometry{margin=2.5cm}

% Configuración para listados de código Python
\definecolor{codegreen}{rgb}{0,0.6,0}
\definecolor{codegray}{rgb}{0.5,0.5,0.5}
\definecolor{codepurple}{rgb}{0.58,0,0.82}
\definecolor{backcolour}{rgb}{0.95,0.95,0.92}

\lstdefinestyle{mystyle}{
    backgroundcolor=\color{backcolour},   
    commentstyle=\color{codegreen},
    keywordstyle=\color{magenta},
    numberstyle=\tiny\color{codegray},
    stringstyle=\color{codepurple},
    basicstyle=\ttfamily\footnotesize,
    breakatwhitespace=false,         
    breaklines=true,                 
    captionpos=b,                    
    keepspaces=true,                 
    numbers=left,                    
    numbersep=5pt,                  
    showspaces=false,                
    showstringspaces=false,
    showtabs=false,                  
    tabsize=4,
    language=Python
}
\lstset{style=mystyle}

\title{\textbf{Análisis Comparativo de Algoritmos de Búsqueda: A* vs Hill Climbing}}
\author{Nombre del Estudiante \\ Ingeniería en Ciencias Informáticas \\ Universidad de las Ciencias Informáticas}
\date{\today}

\begin{document}
\maketitle

\begin{abstract}
Breve resumen del informe (máximo 150 palabras) indicando el problema abordado, los algoritmos comparados y la conclusión principal.
\end{abstract}

\section{Introducción}
Contextualice brevemente el problema de navegación en espacios de estados. Introduzca los algoritmos A* y Hill Climbing, mencionando sus características teóricas (global vs. local, uso de heurísticas).

\section{Metodología}
\subsection{Definición del Espacio de Estados}
Describa cómo se modeló el problema a partir del dataset de coordenadas.
\subsection{Diseño de la Función Heurística}
Explique la función $h(n)$ utilizada y justifique por qué es \textbf{admisible} para el caso de A*.

\section{Resultados Experimentales}
Presente las métricas obtenidas en una tabla comparativa. 
\begin{table}[h!]
\centering
\begin{tabular}{|l|c|c|}
\hline
\textbf{Métrica} & \textbf{A*} & \textbf{Hill Climbing} \\ \hline
¿Llegó al objetivo? &  &  \\ \hline
Nodos expandidos &  &  \\ \hline
Longitud del camino &  &  \\ \hline
Tiempo (seg) &  &  \\ \hline
\end{tabular}
\caption{Comparación de métricas de desempeño.}
\end{table}

Incluya aquí las gráficas generadas con \texttt{matplotlib} mostrando las rutas trazadas por ambos algoritmos.

\section{Análisis y Discusión}
\textbf{Esta es la sección más importante.} Responda críticamente:
\begin{itemize}
    \item ¿Por qué Hill Climbing falló (o tuvo un desempeño inferior) en este escenario específico? Relacione su respuesta con los conceptos de \textit{máximos locales}, \textit{mesetas} o \textit{cresta}.
    \item ¿Cómo garantizó A* la optimalidad del camino? ¿Qué costo computacional pagó por ello?
    \item ¿En qué escenario real usted preferiría usar Hill Climbing sobre A* a pesar de sus limitaciones?
\end{itemize}

\section{Conclusiones}
Síntesis de los aprendizajes obtenidos sobre la idoneidad de cada método según la naturaleza del problema.

\begin{thebibliography}{9}
\bibitem{russell2020} Stuart J. Russell and Peter Norvig. \textit{Artificial Intelligence: A Modern Approach}. 4th Edition. Pearson, 2020.
\end{thebibliography}

\end{document}